\section{Anneaux réguliers}
\label{section:0.17-fr}

\subsection{Définition des anneaux réguliers}
\label{subsection:0.17.1-fr}

\begin{proposition}[17.1.1]
\label{0.17.1.1-fr}
Soient $A$ un anneau local noethérien de dimension $n$, $\mathfrak{m}$ son
idéal maximal, $k=A/\mathfrak{m}$ son corps résiduel. Les conditions suivantes
sont équivalentes~:
\begin{enumerate}
  \item[\rm a)] L'homomorphisme canonique surjectif de $k$-modules gradués
  (15.1.1.1)
  \[
    \varphi:\mathbf{S}^*_k(\mathfrak{m}/\mathfrak{m}^2)
    \longrightarrow\operatorname{gr}^*_{\mathfrak{m}}(A)
  \]
  (où le second membre est le $k$-module gradué associé à $A$ muni de la
  filtration $\mathfrak{m}$-préadique) est bijectif.
  \item[\rm b)] On a
  $\operatorname{rg}_k(\mathfrak{m}/\mathfrak{m}^2)=n$.
  \item[\rm c)] L'idéal $\mathfrak{m}$ admet un système de générateurs de $n$
  éléments.
  \item[\rm d)] L'idéal $\mathfrak{m}$ admet un système de générateurs qui est
  une suite $A$-régulière.
\end{enumerate}
\end{proposition}


\begin{proof}
En vertu du lemme de Nakayama,
$\operatorname{rg}_k(\mathfrak{m}/\mathfrak{m}^2)$ est le plus petit nombre
d'éléments d'un système de générateurs de $\mathfrak{m}$, donc b) et c) sont
équivalentes. D'autre part, si $(x_i)_{1\leq i\leq r}$ est un système de
générateurs de $\mathfrak{m}$ dont les classes $\overline{x}_i$ mod.
$\mathfrak{m}^2$ forment une base du $k$-espace vectoriel
$\mathfrak{m}/\mathfrak{m}^2$, $\mathbf{S}^*(\mathfrak{m}/\mathfrak{m}^2)$
est isomorphe à l'anneau de
\oldpage[0\textsubscript{IV-1}]{40}
polynômes $k[\mathrm{T}_1,\ldots,\mathrm{T}_r]$~; compte tenu de ce que tout
$A$-module de type fini est séparé pour la filtration $\mathfrak{m}$-préadique
($0_{\mathrm I}$, 7.3.5), il résulte de (15.1.9) que les conditions a) et d)
sont équivalentes~; en outre, comme toute suite $A$-régulière d'éléments de
$\mathfrak{m}$ a au plus $n$ éléments (16.4.1), on voit que d) entraîne c).
Reste à prouver que c) entraîne a).

Posons pour abréger
$\mathrm{S}=\mathbf{S}^*(\mathfrak{m}/\mathfrak{m}^2)
=k[\mathrm{T}_1,\ldots,\mathrm{T}_n]$,
$\mathrm{G}=\operatorname{gr}^*_{\mathfrak{m}}(A)$, et considérons la suite
exacte
$0\longrightarrow\mathfrak{J}\longrightarrow\mathrm{S}
\xrightarrow{\varphi}\mathrm{G}\longrightarrow0$, où le noyau
$\mathfrak{J}$ de $\varphi$ est donc un idéal gradué de $\mathrm{S}$. Pour
tout entier $s>0$, on a donc
\begin{equation}
\label{0.17.1.1.1-fr}
\tag{17.1.1.1}
  \binom{s+n-1}{n-1}
  =\operatorname{long}(\mathrm{S}_s)
  =\operatorname{long}(\mathrm{G}_s)
  +\operatorname{long}(\mathfrak{J}_s).
\end{equation}

Supposons $\mathfrak{J}\neq0$, de sorte qu'il existe un élément homogène
$u\in\mathfrak{J}$ de degré $h>0$~; $\mathfrak{J}_s$ contient
$u\mathrm{S}_{s-h}$ pour $s\geq h$, et comme $u\neq0$ et que $\mathrm{S}$ est
intègre, $u\mathrm{S}_{s-h}$ est isomorphe à $\mathrm{S}_{s-h}$ comme
$k$-espace vectoriel~; on aurait donc
$\operatorname{long}(\mathfrak{J}_s)\geq
\operatorname{long}(\mathrm{S}_{s-h})=\binom{s-h+n-1}{n-1}$, d'où pour
$s\geq h$,
\begin{equation}
\label{0.17.1.1.2-fr}
\tag{17.1.1.2}
  \operatorname{long}(\mathrm{G}_s)
  \leq\binom{s+n-1}{n-1}-\binom{s-h+n-1}{n-1}.
\end{equation}
Or, le second membre de (17.1.1.2) est un polynôme en $s$ de degré
$\leq n-2$~; mais on a
$\mathrm{G}_s=\mathfrak{m}^s/\mathfrak{m}^{s+1}$, et pour $s$ assez grand,
$\operatorname{long}(\mathrm{G}_s)$ est un polynôme en $s$ de degré exactement
égal à $n-1$ dont le coefficient dominant est $>0$ (16.2.1), ce qui contredit
l'inégalité (17.1.1.2). C.Q.F.D.
\end{proof}

\begin{definition}[17.1.2]
\label{0.17.1.2-fr}
On dit qu'un anneau local noethérien qui vérifie les conditions équivalentes de
(17.1.1) est \emph{régulier}.
\end{definition}

\begin{corollary}[17.1.3]
\label{0.17.1.3-fr}
Un anneau local régulier est intègre, intégralement clos, et un anneau de
Cohen-Macaulay.
\end{corollary}

\begin{proof}
La dernière assertion résulte de (17.1.1, d))~; d'autre part, si $A$ est
régulier, $\operatorname{gr}^*_{\mathfrak{m}}(A)$ est intègre et complètement
intégralement clos, étant isomorphe à $k[\mathrm{T}_1,\ldots,\mathrm{T}_n]$~;
on en conclut que $A$ possède aussi ces deux propriétés (Bourbaki,
\emph{Alg.~comm.}, chap.~III, \S~2, no~3, cor. de la prop.~1 et chap.~V,
\S~1, no~5, prop.~15).
\end{proof}

\begin{examples}[17.1.4]
\label{0.17.1.4-fr}
\begin{enumerate}
  \item[\rm (i)] Un anneau local régulier de \emph{dimension} $0$, étant
  intègre par (17.1.3), est nécessairement un \emph{corps}, et réciproquement.
  \item[\rm (ii)] Pour qu'un anneau local noethérien de \emph{dimension} $1$
  soit régulier, il faut et il suffit qu'il soit un \emph{anneau de valuation
  discrète}~: en effet, dire qu'un anneau local noethérien de dimension $1$ est
  régulier signifie que son idéal maximal est \emph{principal} et la conclusion
  résulte de Bourbaki, \emph{Alg.~comm.}, chap.~VI, \S~3, no~6, prop.~9.
  \item[\rm (iii)] Soit $k$ un corps~; l'anneau de séries formelles
  $A=k[[\mathrm{T}_1,\ldots,\mathrm{T}_n]]$ est un \emph{anneau régulier de
  dimension $n$}~; en effet, il est clair que les $\mathrm{T}_i$
  $(1\leq i\leq n)$ engendrent l'idéal maximal $\mathfrak{m}$ de $A$ et
  forment une suite $A$-régulière, car
  $A/(\mathrm{T}_1A+\cdots+\mathrm{T}_iA)$ est isomorphe à
  $k[[\mathrm{T}_{i+1},\ldots,\mathrm{T}_n]]$.
\end{enumerate}
\end{examples}

\begin{proposition}[17.1.5]
\label{0.17.1.5-fr}
Pour qu'un anneau local noethérien $A$ soit régulier, il faut et il suffit que
son complété $\widehat{A}$ le soit.
\end{proposition}

\begin{proof}
\oldpage[0\textsubscript{IV-1}]{41}
En effet, l'idéal maximal de $\widehat{A}$ est $\mathfrak{m}\widehat{A}$ et on
sait que $\mathfrak{m}^h/\mathfrak{m}^{h+1}$ et
$(\mathfrak{m}\widehat{A})^h/(\mathfrak{m}\widehat{A})^{h+1}$ sont des
$k$-espaces vectoriels isomorphes~; la condition a) de (17.1.1) est donc la
même pour $A$ et $\widehat{A}$.
\end{proof}

\begin{definition}[17.1.6]
\label{0.17.1.6-fr}
Étant donné un anneau local noethérien $A$, d'idéal maximal $\mathfrak{m}$, on
appelle \emph{système régulier de paramètres de $A$} un système de paramètres
pour $A$ (16.3.5) qui engendre $\mathfrak{m}$.
\end{definition}

L'existence d'un tel système équivaut au fait que $A$ est régulier (17.1.1)~;
si $(x_i)_{1\leq i\leq n}$ est un système régulier de paramètres, c'est un
système minimal de générateurs de $\mathfrak{m}$, dont les classes mod.
$\mathfrak{m}^2$ forment une base de $\mathfrak{m}/\mathfrak{m}^2$ sur $k$~;
en vertu de (17.1.1, a)) et de (15.1.9), $(x_i)$ est une suite $A$-régulière
maximale (d'où la terminologie).

On prendra garde toutefois que, dans un anneau régulier, un système de
paramètres pour $A$ qui est aussi une suite $A$-régulière \emph{n'est pas
nécessairement} un système régulier de paramètres, comme le montre l'exemple
des puissances $k$-ièmes d'un système régulier de paramètres (15.1.20).

\begin{proposition}[17.1.7]
\label{0.17.1.7-fr}
Soient $A$ un anneau local noethérien, $\mathfrak{m}$ son idéal maximal,
$k=A/\mathfrak{m}$ son corps résiduel, $(x_i)_{1\leq i\leq r}$ une suite
d'éléments de $\mathfrak{m}$, $\mathfrak{J}$ l'idéal
$x_1A+\cdots+x_rA$. Les conditions suivantes sont équivalentes~:
\begin{enumerate}
  \item[\rm a)] $A$ est régulier et les $x_i$ font partie d'un système
  régulier de paramètres de $A$.
  \item[\rm a')] $A$ est régulier et les classes $\overline{x}_i$ des $x_i$
  mod. $\mathfrak{m}^2$ sont linéairement indépendantes sur $k$.
  \item[\rm b)] Les $x_i$ font partie d'un système de paramètres pour $A$ et
  $A/\mathfrak{J}$ est un anneau régulier.
\end{enumerate}
En outre, lorsque ces conditions sont vérifiées, $\mathfrak{J}$ est un idéal
premier.
\end{proposition}

\begin{proof}
La dernière assertion résulte trivialement du fait que $A/\mathfrak{J}$,
étant régulier, est intègre (17.1.3). Posons $n=\dim(A)$. L'équivalence de a)
et a') est immédiate~; en effet, les classes mod. $\mathfrak{m}^2$ des éléments
d'un système régulier de paramètres forment une base de
$\mathfrak{m}/\mathfrak{m}^2$ par rapport à $k$, et réciproquement, si les
$\overline{x}_i$ $(1\leq i\leq r)$ sont linéairement indépendants sur $k$, on
peut trouver $n-r$ éléments $x_i\in\mathfrak{m}$ $(r+1\leq i\leq n)$ tels que
les $\overline{x}_i$ $(1\leq i\leq n)$ forment une base de
$\mathfrak{m}/\mathfrak{m}^2$, donc $(x_i)_{1\leq i\leq n}$ est un système
régulier de paramètres. Pour voir que a) entraîne b), remarquons que puisque
les $x_i$ font partie d'un système de paramètres de $A$, on a
$\dim(A/\mathfrak{J})=n-r$ (16.3.6)~; soit
$\mathfrak{n}=\mathfrak{m}/\mathfrak{J}$ l'idéal maximal de
$A/\mathfrak{J}$. On a une suite exacte
\begin{equation}
\label{0.17.1.7.1-fr}
\tag{17.1.7.1}
  0\longrightarrow(\mathfrak{m}^2+\mathfrak{J})/\mathfrak{m}^2
  \longrightarrow\mathfrak{m}/\mathfrak{m}^2
  \longrightarrow\mathfrak{n}/\mathfrak{n}^2\longrightarrow0
\end{equation}
puisque $\mathfrak{n}^2=(\mathfrak{m}^2+\mathfrak{J})/\mathfrak{J}$~;
l'hypothèse a) entraîne que
$\operatorname{rg}_k((\mathfrak{m}^2+\mathfrak{J})/\mathfrak{m}^2)=r$, donc
$\operatorname{rg}_k(\mathfrak{n}/\mathfrak{n}^2)=n-r
=\dim(A/\mathfrak{J})$, ce qui prouve b) en vertu de (17.1.1, b)).

Inversement, pour prouver que b) entraîne a), notons que le fait que les $x_i$
font partie d'un système de paramètres entraîne que
$\dim(A/\mathfrak{J})=n-r$ (16.3.6)~; l'hypothèse que $A/\mathfrak{J}$ est
régulier implique d'autre part
$\operatorname{rg}_k(\mathfrak{n}/\mathfrak{n}^2)=n-r$ (17.1.1)~; par ailleurs
on a évidemment
$\operatorname{rg}_k((\mathfrak{m}^2+\mathfrak{J})/\mathfrak{m}^2)\leq r$,
donc on tire de (17.1.7.1) que
$\operatorname{rg}_k(\mathfrak{m}/\mathfrak{m}^2)\leq r+(n-r)=n$~; donc $A$
est régulier en vertu de (16.2.6) et (17.1.1). En outre la relation
$\operatorname{rg}_k(\mathfrak{m}/\mathfrak{m}^2)=n$ entraîne alors
$\operatorname{rg}_k((\mathfrak{m}^2+\mathfrak{J})/\mathfrak{m}^2)=r$, et par
suite les $\overline{x}_i$ sont linéairement indépendants sur $k$. C.Q.F.D.
\end{proof}

\oldpage[0\textsubscript{IV-1}]{42}

\begin{corollary}[17.1.8]
\label{0.17.1.8-fr}
Soient $A$ un anneau local noethérien, $\mathfrak{m}$ son idéal maximal, $t$ un
élément de $\mathfrak{m}$. Les conditions suivantes sont équivalentes~:
\begin{enumerate}
  \item[\rm a)] $A/tA$ est régulier et $t$ est non diviseur de zéro dans $A$.
  \item[\rm b)] $A$ est régulier et $t\notin\mathfrak{m}^2$.
\end{enumerate}
\end{corollary}

\begin{proof}
C'est le cas particulier $r=1$ de (17.1.7) (compte tenu de ce qu'un élément de
$\mathfrak{m}$ non diviseur de zéro fait partie d'un système de paramètres
(16.4.1)).
\end{proof}

\begin{corollary}[17.1.9]
\label{0.17.1.9-fr}
Soient $A$ un anneau local régulier, $\mathfrak{m}$ son idéal maximal,
$\mathfrak{J}$ un idéal de $A$ contenu dans $\mathfrak{m}$. Les conditions
suivantes sont équivalentes~:
\begin{enumerate}
  \item[\rm a)] L'anneau $A/\mathfrak{J}$ est régulier.
  \item[\rm b)] $\mathfrak{J}$ est engendré par une suite
  $(x_i)_{1\leq i\leq r}$ faisant partie d'un système régulier de paramètres
  de $A$.
\end{enumerate}
\end{corollary}

\begin{proof}
On a déjà vu (17.1.7) que b) entraîne a). Inversement, supposons que
$A/\mathfrak{J}$ soit régulier (ce qui entraîne que $\mathfrak{J}$ est premier)
et soient $n=\dim(A)$, $n-r=\dim(A/\mathfrak{J})$~; avec les notations de
(17.1.7), la suite exacte (17.1.7.1) donne
$\operatorname{rg}_k((\mathfrak{m}^2+\mathfrak{J})/\mathfrak{m}^2)=r$, donc il
existe une suite $(x_i)_{1\leq i\leq r}$ d'éléments de $\mathfrak{J}$ faisant
partie d'un système régulier de paramètres de $A$. Posons
$\mathfrak{J}'=x_1A+\cdots+x_rA$~; il résulte de (17.1.7) que $\mathfrak{J}'$
est un idéal premier de $A$ et que $A/\mathfrak{J}'$ est régulier et de
dimension $n-r$~; mais comme
$A/\mathfrak{J}=(A/\mathfrak{J}')/(\mathfrak{J}/\mathfrak{J}')$ et que
$\mathfrak{J}/\mathfrak{J}'$ est premier dans l'anneau intègre
$A/\mathfrak{J}'$, les dimensions de $A/\mathfrak{J}$ et de
$A/\mathfrak{J}'$ ne peuvent être égales que si
$\mathfrak{J}=\mathfrak{J}'$ (16.1.2.2).
\end{proof}

\subsection{Rappels sur la dimension projective et la dimension injective des
modules}
\label{subsection:0.17.2-fr}

\begin{env}[17.2.1]
\label{0.17.2.1-fr}
Soient $A$ un anneau, $M$ un $A$-module. Rappelons (M, VI, 2) qu'on appelle
\emph{dimension projective} (resp. \emph{dimension injective}) de $M$ et que
l'on note $\operatorname{dim.proj}(M)$ ou
$\operatorname{dim.proj}_A(M)$ (resp. $\operatorname{dim.inj}(M)$ ou
$\operatorname{dim.inj}_A(M)$) le plus petit $n$ (entier ou égal à $+\infty$)
tel qu'il existe une résolution projective gauche (resp. une résolution
injective droite) de $M$ qui soit \emph{de longueur $n$}. Il revient au même
(\emph{loc. cit.}) de dire que $n$ est le plus petit nombre tel que pour tout
$A$-module $N$, on ait $\operatorname{Ext}^i_A(M,N)=0$ (resp.
$\operatorname{Ext}^i_A(N,M)=0$) pour tout $i>n$, ou seulement pour $i=n+1$.
Pour tout \emph{facteur direct} $M'$ de $M$, on a donc
$\operatorname{dim.proj}(M')\leq\operatorname{dim.proj}(M)$ puisque
$\operatorname{Ext}^i_A(M',N)$ est facteur direct de
$\operatorname{Ext}^i_A(M,N)$~; de même
$\operatorname{dim.inj}(M')\leq\operatorname{dim.inj}(M)$. Une condition
équivalente à $\operatorname{dim.proj}(M)=n$ (resp.
$\operatorname{dim.inj}(M)=n$) est que $n$ est le plus petit nombre tel que,
pour toute suite exacte
\[
  0\longrightarrow R\longrightarrow P_{n-1}\longrightarrow\cdots
  \longrightarrow P_0\longrightarrow M\longrightarrow0
\]
(resp.
$0\longrightarrow M\longrightarrow Q_0\longrightarrow\cdots
\longrightarrow Q_{n-1}\longrightarrow C\longrightarrow0$), où tous les
$P_i$ pour $0\leq i<n$ sont projectifs (resp. tous les $Q_i$ pour $0\leq i<n$
sont injectifs), alors $R$ est projectif (resp. $C$ est injectif).
\end{env}

\begin{remarks}[17.2.2]
\label{0.17.2.2-fr}
\begin{enumerate}
  \item[\rm (i)] Dire que $M$ est un $A$-module \emph{projectif} (resp.
  \emph{injectif}) équivaut donc à dire que
  $\operatorname{dim.proj}(M)=0$ (resp. $\operatorname{dim.inj}(M)=0$).
  \item[\rm (ii)] Soit $T$ un \emph{foncteur covariant} additif de la catégorie
  des $A$-modules dans une catégorie abélienne. Il résulte aussitôt des
  définitions de (17.2.1) et de celle des
  \oldpage[0\textsubscript{IV-1}]{43}
  foncteurs dérivés que si $\operatorname{dim.proj}(M)\leq n$ (resp.
  $\operatorname{dim.inj}(M)\leq n$), on a $L^iT(M)=0$ (resp.
  $R^iT(M)=0$) \emph{pour tout $i>n$}.
  \item[\rm (iii)] Si on suppose $A$ \emph{noethérien} et $M$ \emph{de type
  fini}, la dernière interprétation de la dimension projective donnée dans
  (17.2.1) montre que si $\operatorname{dim.proj}(M)=n$, $M$ admet une
  résolution gauche de longueur $n$ formée de modules projectifs \emph{de type
  fini}. Si de plus $A$ est un anneau \emph{local} noethérien, $M$ admet une
  résolution de longueur $n$ par des modules \emph{libres de type fini}, un
  $A$-module projectif de type fini étant alors libre (Bourbaki,
  \emph{Alg.~comm.}, chap.~II, \S~3, no~2, cor.~2 de la prop.~5).
\end{enumerate}
\end{remarks}

\begin{lemma}[17.2.3]
\label{0.17.2.3-fr}
Soient $A$ un anneau, $M$ un $A$-module. Pour que
$\operatorname{dim.inj}(M)\leq n$, il faut et il suffit que pour tout
$A$-module monogène $N$, on ait $\operatorname{Ext}^{n+1}_A(N,M)=0$.
\end{lemma}

\begin{proof}
Avec les notations de (17.2.1), il s'agit de prouver que $C$ est injectif~;
pour tout $A$-module $N$, $\operatorname{Ext}^{n+1}_A(N,M)$ est isomorphe à
$\operatorname{Ext}^1_A(N,C)$ (M, V, 7), donc on a
$\operatorname{Ext}^1_A(N,C)=0$ pour tout $A$-module $N$ de la forme
$A/\mathfrak{J}$, où $\mathfrak{J}$ est un idéal quelconque de $A$. La suite
exacte des Ext montre alors que l'homomorphisme canonique
$\operatorname{Hom}(A,C)\longrightarrow\operatorname{Hom}(\mathfrak{J},C)$ est
surjectif pour tout idéal $\mathfrak{J}$ de $A$, ce qui implique que $C$ est un
$A$-module injectif (M, I, 3.2).
\end{proof}

\begin{lemma}[17.2.4]
\label{0.17.2.4-fr}
Soient $A$ un anneau noethérien, $M$ un $A$-module de type fini. Pour que
$\operatorname{dim.proj}(M)\leq n$, il faut et il suffit que pour tout
$A$-module monogène $N$, on ait $\operatorname{Ext}^{n+1}_A(M,N)=0$.
\end{lemma}

\begin{proof}
On sait en effet (M, VI, 2.5) que la condition
$\operatorname{dim.proj}(M)\leq n$ est équivalente à
$\operatorname{Ext}^{n+1}_A(M,N)=0$ pour tout $A$-module $N$ \emph{de type
fini}~; pour voir que la condition de l'énoncé est aussi suffisante, on
raisonne par récurrence sur le nombre de générateurs $m$ de $N$~: il y a un
sous-module $N_1$ de $N$ engendré par $m-1$ éléments et tel que
$N_2=N/N_1$ soit monogène~; de la suite exacte
$0\longrightarrow N_1\longrightarrow N\longrightarrow N_2\longrightarrow0$,
on déduit alors la suite exacte
$\operatorname{Ext}^{n+1}_A(M,N_1)\longrightarrow
\operatorname{Ext}^{n+1}_A(M,N)\longrightarrow
\operatorname{Ext}^{n+1}_A(M,N_2)$ et l'hypothèse de récurrence montre que la
condition de l'énoncé entraîne bien $\operatorname{Ext}^{n+1}_A(M,N)=0$.
\end{proof}

\begin{corollary}[17.2.5]
\label{0.17.2.5-fr}
Soient $A$ un anneau noethérien, $M$ un $A$-module. On a
\begin{equation}
\label{0.17.2.5.1-fr}
\tag{17.2.5.1}
  \operatorname{dim.inj}_A(M)
  =\sup_{\mathfrak{m}}
  \bigl(\operatorname{dim.inj}_{A_{\mathfrak{m}}}(M_{\mathfrak{m}})\bigr)
\end{equation}
et si $M$ est \emph{de type fini}
\begin{equation}
\label{0.17.2.5.2-fr}
\tag{17.2.5.2}
  \operatorname{dim.proj}_A(M)
  =\sup_{\mathfrak{m}}
  \bigl(\operatorname{dim.proj}_{A_{\mathfrak{m}}}(M_{\mathfrak{m}})\bigr)
\end{equation}
où $\mathfrak{m}$ parcourt l'ensemble des idéaux premiers (ou l'ensemble des
idéaux maximaux) de $A$.
\end{corollary}

\begin{proof}
En effet, si $N$ est un $A$-module de type fini (donc de présentation finie),
on a, pour toute partie multiplicative $S$ de $A$ et pour tout $i\geq0$
\[
  S^{-1}\operatorname{Ext}^i_A(N,M)
  =\operatorname{Ext}^i_{S^{-1}A}(S^{-1}N,S^{-1}M)
\]
par platitude, en considérant une résolution libre de $M$ et en utilisant le
fait que la relation précédente est vraie pour $i=0$ (Bourbaki,
\emph{Alg.~comm.}, chap.~II, \S~2, no~7, prop.~19). En particulier
$\operatorname{Ext}^i_{A_{\mathfrak{m}}}
(A_{\mathfrak{m}}/\mathfrak{J}A_{\mathfrak{m}},M_{\mathfrak{m}})
=(\operatorname{Ext}^i_A(A/\mathfrak{J},M))_{\mathfrak{m}}$ pour tout idéal
premier $\mathfrak{m}$ et tout idéal $\mathfrak{J}$ de $A$~; si l'on tient
compte de (17.2.3) et de ce que tout
\oldpage[0\textsubscript{IV-1}]{44}
idéal de $A_{\mathfrak{m}}$ est de la forme $\mathfrak{J}A_{\mathfrak{m}}$
pour un idéal convenable $\mathfrak{J}$ de $A$, on déduit la formule
(17.2.5.1) de ce qui précède et de Bourbaki, \emph{Alg.~comm.}, chap.~II,
\S~3, no~3, th.~1. On procède de même pour (17.2.5.2) en utilisant cette fois
(17.2.4) et en échangeant les rôles de $M$ et de $N$.
\end{proof}

Pour les anneaux noethériens et les modules de type fini, l'étude de la
dimension projective ou de la dimension injective est donc ramenée par
(17.2.5) au cas des anneaux \emph{locaux}. On a alors le

\begin{lemma}[17.2.6]
\label{0.17.2.6-fr}
Soient $A$ un anneau local noethérien, $k$ son corps résiduel, $M$ un
$A$-module de type fini. Pour que $\operatorname{dim.proj}(M)\leq n$, il est
nécessaire que $\operatorname{Tor}^A_i(M,k)=0$ pour $i>n$, et il suffit que
$\operatorname{Tor}^A_{n+1}(M,k)=0$.
\end{lemma}

\begin{proof}
La nécessité est un cas particulier de la remarque (17.2.2 (ii)), appliquée au
foncteur covariant $M\mapsto k\otimes_A M$. Pour prouver que la condition est
suffisante, il s'agit, avec les notations de (17.2.1), d'établir que $R$ est
projectif lorsque les $P_i$ sont supposés de type fini~; or
$\operatorname{Tor}^A_{n+1}(M,k)$ est isomorphe à
$\operatorname{Tor}^A_1(R,k)$ (M, V, 7)~; et on sait que, puisque $R$ est de
type fini, la condition $\operatorname{Tor}^A_1(R,k)=0$ entraîne que $R$ est
libre (Bourbaki, \emph{Alg.~comm.}, chap.~II, \S~3, no~2, cor.~2 de la
prop.~5).
\end{proof}

\begin{corollary}[17.2.7]
\label{0.17.2.7-fr}
Sous les hypothèses de (17.2.6), soit $x$ un élément $M$-régulier de $A$
appartenant à l'idéal maximal de $A$~; alors on a
\begin{equation}
\label{0.17.2.7.1-fr}
\tag{17.2.7.1}
  \operatorname{dim.proj}(M/xM)=\operatorname{dim.proj}(M)+1.
\end{equation}
\end{corollary}

\begin{proof}
En effet, de la suite exacte
$0\longrightarrow M\xrightarrow{x}M\longrightarrow M/xM\longrightarrow0$
définie par la multiplication par $x$ dans $M$, on déduit la suite exacte des
Tor~:
\[
  \operatorname{Tor}^A_i(M,k)\xrightarrow{x}
  \operatorname{Tor}^A_i(M,k)\longrightarrow
  \operatorname{Tor}^A_i(M/xM,k)\longrightarrow
  \operatorname{Tor}^A_{i-1}(M,k)\xrightarrow{x}
  \operatorname{Tor}^A_{i-1}(M,k)
\]
pour tout $i\geq1$. Or, pour tout $A$-module $N$, l'homothétie de rapport $x$
dans $M\otimes_A N$ provient aussi bien de l'homothétie de rapport $x$ dans
$M$ que de l'homothétie de rapport $x$ dans $N$~; on en conclut aussitôt,
puisque l'homothétie de rapport $x$ dans $k$ est nulle par définition, que,
pour tout $i$, l'homomorphisme
$\operatorname{Tor}^A_i(M,k)\xrightarrow{x}\operatorname{Tor}^A_i(M,k)$ est
\emph{nul}~; autrement dit, on a la suite exacte
\[
  0\longrightarrow\operatorname{Tor}^A_i(M,k)\longrightarrow
  \operatorname{Tor}^A_i(M/xM,k)\longrightarrow
  \operatorname{Tor}^A_{i-1}(M,k)\longrightarrow0.
\]
Si $\operatorname{dim.proj}(M)=n$, on a
$\operatorname{Tor}^A_n(M,k)\neq0$ et
$\operatorname{Tor}^A_{n+1}(M,k)=\operatorname{Tor}^A_{n+2}(M,k)=0$ en vertu
de (17.2.6). Il résulte donc de ce qui précède que l'on a
$\operatorname{Tor}^A_{n+1}(M/xM,k)\neq0$ et
$\operatorname{Tor}^A_{n+2}(M/xM,k)=0$, donc
$\operatorname{dim.proj}(M/xM)=n+1$ par (17.2.6).
\end{proof}

\begin{proposition}[17.2.8]
\label{0.17.2.8-fr}
\textnormal{(M.~Auslander)} Soit $A$ un anneau. Les conditions suivantes sont
équivalentes~:
\begin{enumerate}
  \item[\rm a)] Tout $A$-module est de dimension projective $\leq n$.
  \item[\rm a')] Tout $A$-module de type fini est de dimension projective
  $\leq n$.
  \item[\rm b)] Tout $A$-module est de dimension injective $\leq n$.
  \item[\rm c)] Pour tout couple de $A$-modules $M$, $N$ on a
  $\operatorname{Ext}^{n+1}_A(M,N)=0$.
  \item[\rm c')] Pour tout couple de $A$-modules $M$, $N$ tel que $M$ soit de
  type fini (ou seulement monogène), on a
  $\operatorname{Ext}^{n+1}_A(M,N)=0$.
\end{enumerate}
\end{proposition}

\begin{proof}
\oldpage[0\textsubscript{IV-1}]{45}
Cela résulte aussitôt de (17.2.1) et (17.2.3).
\end{proof}

Le plus petit nombre $n$ (entier ou $+\infty$) pour lequel les conditions
équivalentes de (17.2.8) sont satisfaites est appelé \emph{dimension
cohomologique globale} (ou simplement \emph{dimension cohomologique}) de $A$
et noté $\operatorname{dim.coh}(A)$.

\begin{proposition}[17.2.9]
\label{0.17.2.9-fr}
Soit $A$ un anneau noethérien. Les conditions suivantes sont équivalentes~:
\begin{enumerate}
  \item[\rm a)] $\operatorname{dim.coh}(A)\leq n$.
  \item[\rm b)] Tout $A$-module de type fini est de dimension injective
  $\leq n$.
  \item[\rm c)] Pour tout couple de $A$-modules de type fini, $M$, $N$, on a
  $\operatorname{Ext}^{n+1}_A(M,N)=0$.
\end{enumerate}
\end{proposition}

\begin{proof}
Cela résulte aussitôt de la définition et de (17.2.4).
\end{proof}

\begin{corollary}[17.2.10]
\label{0.17.2.10-fr}
Si $A$ est un anneau noethérien, on a
\begin{equation}
\label{0.17.2.10.1-fr}
\tag{17.2.10.1}
  \operatorname{dim.coh}(A)
  =\sup_{\mathfrak{m}}
  \bigl(\operatorname{dim.coh}(A_{\mathfrak{m}})\bigr)
\end{equation}
où $\mathfrak{m}$ parcourt le spectre de $A$ (ou l'ensemble des idéaux
maximaux de $A$).
\end{corollary}

\begin{proof}
Cela résulte de (17.2.9) et (17.2.5).
\end{proof}

\begin{proposition}[17.2.11]
\label{0.17.2.11-fr}
Soient $A$ un anneau local noethérien, $k$ son corps résiduel. Pour que
$\operatorname{dim.coh}(A)\leq n$, il est nécessaire que
$\operatorname{Tor}^A_i(k,k)=0$, pour $i>n$, et il suffit que
$\operatorname{Tor}^A_{n+1}(k,k)=0$.
\end{proposition}

\begin{proof}
Compte tenu de (17.2.6), il suffit de prouver que les relations
$\operatorname{dim.coh}(A)\leq n$ et
$\operatorname{dim.proj}_A(k)\leq n$ sont équivalentes. Il est clair que la
première entraîne la seconde par définition. Inversement, si
$\operatorname{dim.proj}_A(k)\leq n$, on a
$\operatorname{Tor}^A_i(M,k)=0$ pour $i>n$ et pour tout $A$-module $M$ en
vertu de (17.2.2, (ii))~; donc $\operatorname{dim.proj}(M)\leq n$, ce qui
prouve la proposition en vertu de (17.2.8).
\end{proof}

\begin{corollary}[17.2.12]
\label{0.17.2.12-fr}
Soit $A$ un anneau noethérien. Pour que l'on ait
$\operatorname{dim.coh}(A)\leq n$, il est nécessaire que, pour tout idéal
maximal $\mathfrak{m}$ de $A$, on ait
$\operatorname{Tor}^{A_{\mathfrak{m}}}_i(A/\mathfrak{m},A/\mathfrak{m})=0$
pour $i>n$, et il suffit que ces relations soient vérifiées pour $i=n+1$.
\end{corollary}

\begin{proof}
Cela résulte aussitôt de (17.2.11) et (17.2.10).
\end{proof}

\begin{proposition}[17.2.13]
\label{0.17.2.13-fr}
Soient $A$, $B$ deux anneaux locaux noethériens,
$\varphi:A\longrightarrow B$ un homomorphisme local faisant de $B$ un
$A$-module plat. Alors on a
\begin{equation}
\label{0.17.2.13.1-fr}
\tag{17.2.13.1}
  \operatorname{dim.coh}(A)\leq\operatorname{dim.coh}(B).
\end{equation}
\end{proposition}

\begin{proof}
Supposons que $\operatorname{dim.coh}(B)=n$ soit finie~; il suffit de prouver
que pour tout couple $(M,N)$ de $A$-modules de type fini, on a
$\operatorname{Tor}^A_i(M,N)=0$ pour $i>n$ (17.2.11). Comme $B$ est un
$A$-module fidèlement plat ($0_{\rm I}$, 6.6.2), il revient au même
($0_{\rm I}$, 6.4.1) de montrer que l'on a
\begin{equation}
\label{0.17.2.13.2-fr}
\tag{17.2.13.2}
  \operatorname{Tor}^A_i(M,N)\otimes_A B=0.
\end{equation}
Or, si $L_\bullet=(L_j)$ est une résolution droite de $M$ par des $A$-modules
libres, il résulte de ce que $B$ est un $A$-module plat que
$L_\bullet\otimes_A B=(L_j\otimes_A B)$ est une résolution droite du
$B$-module $M\otimes_A B$ par des $B$-modules libres~; en outre, on a
$(L_j\otimes_A B)\otimes_B(N\otimes_A B)
=(L_j\otimes_A N)\otimes_A B$, d'où on conclut aussitôt que le premier membre
\oldpage[0\textsubscript{IV-1}]{46}
de (17.2.13.2) est égal à
$\operatorname{Tor}^B_i(M\otimes_A B,N\otimes_A B)$~; l'hypothèse sur $B$
entraîne que ce $B$-module est nul pour $i>n$ (17.2.2, (ii)), d'où la
conclusion.
\end{proof}

\begin{env}[17.2.14]
\label{0.17.2.14-fr}
Soient $(X,\mathscr{O}_X)$ un espace annelé, $\mathscr{F}$ un
$\mathscr{O}_X$-Module~; on appelle \emph{dimension projective} (resp.
\emph{injective}) \emph{ponctuelle} de $\mathscr{F}$ et on note
$\operatorname{dim.proj}(\mathscr{F})$ (resp.
$\operatorname{dim.inj}(\mathscr{F})$) le nombre
$\sup_{x\in X}(\operatorname{dim.proj}(\mathscr{F}_x))$ (resp.
$\sup_{x\in X}(\operatorname{dim.inj}(\mathscr{F}_x))$. On appelle
\emph{dimension cohomologique ponctuelle de $X$} et on note
$\operatorname{dim.coh}(X)$ le nombre
$\sup_{x\in X}(\operatorname{dim.coh}(\mathscr{O}_x))$. Il résulte de
(17.2.5) et (17.2.10) que lorsque $A$ est un anneau noethérien, $M$ un
$A$-module de type fini et $X=\operatorname{Spec}(A)$, on a
$\operatorname{dim.proj}(\widetilde{M})=\operatorname{dim.proj}(M)$,
$\operatorname{dim.inj}(\widetilde{M})=\operatorname{dim.inj}(M)$ et
$\operatorname{dim.coh}(X)=\operatorname{dim.coh}(A)$. On appelle
\emph{dimension projective} (resp. \emph{injective}) \emph{de $\mathscr{F}$ en
un point $x\in X$} la dimension projective (resp. injective) de
$\mathscr{F}_x$, \emph{dimension cohomologique de $X$ en un point $x$} la
dimension cohomologique de $\mathscr{O}_x$.
\end{env}

\begin{proposition}[17.2.15]
\label{0.17.2.15-fr}
Soient $X$, $Y$ deux espaces annelés en anneaux locaux noethériens,
$f:X\longrightarrow Y$ un morphisme plat. Si
$\operatorname{dim.coh}(X)\leq n$, alors $Y$ est de dimension cohomologique
$\leq n$ en tout point de $f(X)$.
\end{proposition}

\begin{proof}
Cela résulte aussitôt de (17.2.13).
\end{proof}

\subsection{Théorie cohomologique des anneaux réguliers}
\label{subsection:0.17.3-fr}

\begin{theorem}[17.3.1]
\label{0.17.3.1-fr}
\textnormal{(Hilbert-Serre)} Soit $A$ un anneau local noethérien. Pour que $A$
soit de dimension cohomologique finie, il faut et il suffit que $A$ soit
régulier~; on a alors
\begin{equation}
\label{0.17.3.1.1-fr}
\tag{17.3.1.1}
  \operatorname{dim.coh}(A)=\dim(A).
\end{equation}
\end{theorem}

\begin{proof}
Supposons $A$ régulier~; soient $\mathfrak{m}$ son idéal maximal,
$\mathbf{x}=(x_i)_{1\leq i\leq n}$ un système régulier de paramètres de $A$
(17.1.6)~; considérons le complexe de l'algèbre extérieure
$K_\bullet(\mathbf{x})$ (\textbf{III}, 1.1.1), qui est formé de $A$-modules
libres, avec $K_i(\mathbf{x})=0$ pour $i>n$~; comme $(x_i)$ est une suite
régulière, on a $H_i(K_\bullet(\mathbf{x}))=0$ pour $i>0$ (\textbf{III},
1.1.4 et 1.1.3.3) et
$H_0(K_\bullet(\mathbf{x}))=A/(x_1A+\cdots+x_nA)=A/\mathfrak{m}=k$~; les
$K_i(\mathbf{x})$ forment par suite une \emph{résolution libre} de $k$ de
longueur $n$. Or, le fait que les $x_i$ appartiennent à $\mathfrak{m}$
entraîne aussitôt que dans le complexe
$K_\bullet(\mathbf{x},k)=K_\bullet(\mathbf{x})\otimes_A k$, l'opérateur bord
est nul en toutes dimensions, si bien que l'on a, par définition,
$\operatorname{Tor}^A_i(k,k)=\bigwedge^i(k^n)$~; l'égalité (17.3.1.1) résulte
donc aussitôt de (17.2.11) (ce résultat est essentiellement le «~théorème des
syzygies~» de Hilbert).

Montrons maintenant que si $A$ est un anneau local noethérien,
$\mathfrak{m}$ son idéal maximal, et si
$\operatorname{dim.proj}_A(\mathfrak{m})$ est \emph{finie}, alors $A$ est
régulier, ce qui achèvera de démontrer (17.3.1). Nous procéderons par
récurrence sur $n=\operatorname{rg}_k(\mathfrak{m}/\mathfrak{m}^2)$. Pour
$n=0$, on a $\mathfrak{m}=0$ et l'assertion est triviale.
\end{proof}

\begin{lemma}[17.3.1.2]
\label{0.17.3.1.2-fr}
\textnormal{(Nagata)} Si tout élément de
$\mathfrak{m}-\mathfrak{m}^2$ est diviseur de zéro dans $A$, il existe
$c\neq0$ dans $A$ tel que $c\mathfrak{m}=0$ (autrement dit on a
$\mathfrak{m}\in\operatorname{Ass}(A)$).
\end{lemma}

\begin{proof}
On peut se borner au cas où $\mathfrak{m}\neq0$, donc
$\mathfrak{m}\neq\mathfrak{m}^2$. L'hypothèse entraîne que
$\mathfrak{m}-\mathfrak{m}^2$ est contenu dans la réunion des idéaux
$\mathfrak{p}_i$ de $\operatorname{Ass}(A)$ (Bourbaki,
\emph{Alg.~comm.}, chap.~IV, \S~1, no~1, cor.~3 de la prop.~2)~; donc
$\mathfrak{m}$ est contenu dans la réunion de $\mathfrak{m}^2$
\oldpage[0\textsubscript{IV-1}]{47}
et des $\mathfrak{p}_i$, donc dans l'un des $\mathfrak{p}_i$, puisque
$\mathfrak{m}\not\subset\mathfrak{m}^2$ (Bourbaki, \emph{Alg.~comm.},
chap.~II, \S~1, no~1, prop.~2)~; $\mathfrak{m}$ étant maximal, cela prouve le
lemme.
\end{proof}

\begin{lemma}[17.3.1.3]
\label{0.17.3.1.3-fr}
Si $a\in\mathfrak{m}-\mathfrak{m}^2$, $\mathfrak{m}/Aa$ est isomorphe à un
facteur direct de $\mathfrak{m}/a\mathfrak{m}$.
\end{lemma}

\begin{proof}
Il existe un système minimal de générateurs de $\mathfrak{m}$ contenant $a$
(dont les images dans $\mathfrak{m}/\mathfrak{m}^2$ forment une base de ce
$k$-espace vectoriel)~; soit $\mathfrak{b}$ l'idéal engendré par les éléments
de ce système autres que $a$. Comme la relation $xa\in\mathfrak{b}$ entraîne
$x\in\mathfrak{m}$ (en considérant l'image de $xa$ dans
$\mathfrak{m}/\mathfrak{m}^2$), on a
$\mathfrak{b}\cap Aa\subset a\mathfrak{m}$ d'où un homomorphisme
$\mathfrak{b}/(\mathfrak{b}\cap Aa)\longrightarrow
\mathfrak{m}/a\mathfrak{m}$ déduit de l'injection
$\mathfrak{b}\longrightarrow\mathfrak{m}$ par passage aux quotients, et qui
est \emph{injectif}. En outre, $\mathfrak{b}+Aa=\mathfrak{m}$, et
l'homomorphisme composé
\[
  \mathfrak{m}/Aa=(\mathfrak{b}+Aa)/Aa
  \xrightarrow{\sim}\mathfrak{b}/(\mathfrak{b}\cap Aa)
  \longrightarrow\mathfrak{m}/a\mathfrak{m}
  \longrightarrow\mathfrak{m}/Aa
\]
est l'identité~; d'où le lemme.
\end{proof}

\begin{lemma}[17.3.1.4]
\label{0.17.3.1.4-fr}
Soient $A$ un anneau local noethérien, $\mathfrak{m}$ son idéal maximal, $E$
un $A$-module de type fini et de dimension projective finie. Si
$a\in\mathfrak{m}$ est $A$-régulier et $E$-régulier, alors $E/aE$ est un
$(A/aA)$-module de dimension projective finie, au plus égale à
$\operatorname{dim.proj}_A(E)$.
\end{lemma}

\begin{proof}
Raisonnons par récurrence sur $h=\operatorname{dim.proj}_A(E)$, le cas $h=0$
étant trivial puisque alors $E$ est un $A$-module projectif, donc $E/aE$ un
$(A/aA)$-module projectif. Il existe une suite exacte
\[
  0\longrightarrow N\longrightarrow L\longrightarrow E\longrightarrow0
\]
où $L$ est libre et $\operatorname{dim.proj}_A(N)=h-1$ (17.2.2, (iii)), $N$
étant de type fini. En outre, la suite
\[
  0\longrightarrow N/aN\longrightarrow L/aL\longrightarrow E/aE
  \longrightarrow0
\]
est exacte (15.1.18). Comme $a$ est $A$-régulier, il est aussi $N$-régulier
puisque $L$ est libre~; l'hypothèse de récurrence entraîne que $N/aN$ est un
$(A/aA)$-module de dimension projective $\leq h-1$, et comme $L/aL$ est un
$(A/aA)$-module libre, $E/aE$ est un $(A/aA)$-module de dimension projective
$\leq h$.
\end{proof}

Examinons maintenant deux cas~:

\medskip\noindent\textbf{I.} --- Supposons d'abord que tout élément de
$\mathfrak{m}-\mathfrak{m}^2$ soit diviseur de zéro dans $A$, auquel cas
(17.3.1.2) il existe $c\neq0$ dans $A$ tel que $c\mathfrak{m}=0$. Montrons
alors que $\mathfrak{m}=0$. S'il n'en était pas ainsi, notons d'abord que
$\mathfrak{m}$ ne pourrait être un $A$-module projectif, car il serait libre
(Bourbaki, \emph{Alg.~comm.}, chap.~II, \S~3, no~2, cor.~2 de la prop.~5),
ce qui contredit la relation $c\mathfrak{m}=0$. On aurait donc
$n=\operatorname{dim.coh}(A)\geq1$. Comme
$\mathfrak{m}\in\operatorname{Ass}(A)$, il existerait une suite exacte de
$A$-homomorphismes
\[
  0\longrightarrow k\longrightarrow A\longrightarrow E\longrightarrow0.
\]
Mais cela est absurde, car en vertu de la relation $n\geq1$, la suite exacte
des Tor donnerait la suite exacte
$0\longrightarrow\operatorname{Tor}^A_{n+1}(E,k)\longrightarrow
\operatorname{Tor}^A_n(k,k)\longrightarrow0$~; or on a
$\operatorname{Tor}^A_{n+1}(E,k)=0$ (17.2.2, (ii)) et
$\operatorname{Tor}^A_n(k,k)\neq0$ (17.2.11) et nous avons abouti à une
contradiction.

\medskip\noindent\textbf{II.} --- On peut donc se borner au cas où il existe
$a\in\mathfrak{m}-\mathfrak{m}^2$ qui est un élément $A$-régulier, et par
suite aussi $\mathfrak{m}$-régulier. Considérons l'anneau $A'=A/aA$ et son
idéal
\oldpage[0\textsubscript{IV-1}]{48}
maximal $\mathfrak{m}'=\mathfrak{m}/aA$~; il est clair que
$\operatorname{rg}_k(\mathfrak{m}'/\mathfrak{m}'^2)=n-1$. En vertu de
(17.3.1.4), $\mathfrak{m}/a\mathfrak{m}$ est un $A'$-module de dimension
projective finie, donc il en est de même de
$\mathfrak{m}'=\mathfrak{m}/aA$, qui en est un facteur direct (17.3.1.3 et
17.2.1). L'hypothèse de récurrence entraîne par suite que $A'=A/aA$ est
régulier, ce qui prouve par (17.1.8) que $A$ est régulier.

\begin{corollary}[17.3.2]
\label{0.17.3.2-fr}
Si $A$ est un anneau local régulier, $A_{\mathfrak{p}}$ est régulier pour tout
idéal premier $\mathfrak{p}$ de $A$.
\end{corollary}

\begin{proof}
En effet, on a vu (17.2.10) que
$\operatorname{dim.coh}(A_{\mathfrak{p}})\leq
\operatorname{dim.coh}(A)$, donc la conclusion résulte aussitôt de (17.3.1).
\end{proof}

\begin{proposition}[17.3.3]
\label{0.17.3.3-fr}
Soient $A$, $B$ deux anneaux locaux noethériens,
$\varphi:A\longrightarrow B$ un homomorphisme local, $\mathfrak{m}$ (resp.
$\mathfrak{n}$) l'idéal maximal de $A$ (resp. $B$),
$k=A/\mathfrak{m}$ (resp. $k'=B/\mathfrak{n}$) le corps résiduel de $A$
(resp. $B$). On a donc un homomorphisme canonique de $k'$-espaces vectoriels
\begin{equation}
\label{0.17.3.3.1-fr}
\tag{17.3.3.1}
  \psi:(\mathfrak{m}/\mathfrak{m}^2)\otimes_k k'
  \longrightarrow\mathfrak{n}/\mathfrak{n}^2.
\end{equation}
\begin{enumerate}
  \item[\rm (i)] Si $B$ est régulier et si $B$ est un $A$-module plat, $A$
  est régulier.
  \item[\rm (ii)] Les conditions suivantes sont équivalentes~:
  \begin{enumerate}
    \item[\rm a)] $B$ est régulier et l'homomorphisme $\psi$ (17.3.3.1) est
    injectif.
    \item[\rm b)] $A$ et $B$ sont réguliers, et pour un système régulier de
    paramètres $(x_i)_{1\leq i\leq m}$ de $A$, les $\varphi(x_i)$ font partie
    d'un système régulier de paramètres de $B$ (auquel cas cette propriété a
    lieu pour \emph{tout} système régulier de paramètres de $A$).
    \item[\rm c)] $B$ et $B\otimes_A k$ sont réguliers, et $B$ est un
    $A$-module plat.
    \item[\rm d)] $A$ et $B\otimes_A k$ sont réguliers, et $B$ est un
    $A$-module plat.
    \item[\rm e)] $A$ et $B\otimes_A k$ sont réguliers, et l'on a
    \begin{equation}
    \label{0.17.3.3.2-fr}
    \tag{17.3.3.2}
      \dim(B)=\dim(A)+\dim(B\otimes_A k).
    \end{equation}
  \end{enumerate}
\end{enumerate}
\end{proposition}

\begin{proof}
\textnormal{(i)} On a
$\operatorname{dim.coh}(A)\leq\operatorname{dim.coh}(B)$ par (17.2.13), donc
il suffit d'appliquer (17.3.1).

\textnormal{(ii)} Lorsque $A$ est supposé régulier et que $(x_i)$ est un
système régulier de paramètres de $A$, dire que $B$ est un $A$-module
\emph{plat} équivaut, en vertu de (15.1.21) à dire que la suite des
$y_i=\varphi(x_i)$ est $B$-régulière (puisque $B\otimes_A k$ est un
$k$-module plat). D'autre part, comme $\dim(A)=m$, et que la suite $(y_i)$
engendre l'idéal $\mathfrak{m}B$, la relation (17.3.3.2), qui s'écrit aussi
$\dim(B/\mathfrak{m}B)=\dim(B)-m$, équivaut à dire que la suite
$(y_i)_{1\leq i\leq m}$ fait partie d'un système de paramètres de $B$
(16.3.7). On voit donc que les conditions b), d) et e) équivalent
respectivement aux suivantes~:
\begin{enumerate}
  \item[\rm b')] $A$ et $B$ sont réguliers et $(y_i)_{1\leq i\leq m}$ fait
  partie d'un système régulier de paramètres de $B$.
  \item[\rm d')] $A$ est régulier,
  $B/(\sum_{i=1}^m y_iB)$ est régulier et la suite $(y_i)$ est
  $B$-régulière.
  \item[\rm e')] $A$ est régulier,
  $B/(\sum_{i=1}^m y_iB)$ est régulier et la suite $(y_i)$ fait partie d'un
  système de paramètres de $B$.
\end{enumerate}
Or, b') et e') sont équivalentes en vertu de (17.1.7), et comme d') entraîne
e') (16.4.1) et est entraînée par b') (17.1.7), elle leur est équivalente. La
conjonction
\oldpage[0\textsubscript{IV-1}]{49}
de b) et d) implique trivialement c), et en vertu de (i), c) entraîne d)~; on
a donc prouvé l'équivalence de b), c), d) et e). En outre, il est clair que b)
entraîne a), les classes des $x_i$ dans
$\mathfrak{m}/\mathfrak{m}^2$ formant alors une base de ce $k$-espace
vectoriel et les classes des $y_i$ dans $\mathfrak{n}/\mathfrak{n}^2$ un
système libre de ce $k'$-espace vectoriel. Reste donc à prouver que a)
entraîne e). Posons $V=\mathfrak{m}/\mathfrak{m}^2$,
$W=\mathfrak{n}/\mathfrak{n}^2$~; il est immédiat que l'on a
$\psi(V\otimes_k k')=(\mathfrak{n}^2+\mathfrak{m}B)/\mathfrak{n}^2$. On a
d'une part, en vertu de (16.3.9)
\begin{equation}
\label{0.17.3.3.3-fr}
\tag{17.3.3.3}
  \dim(B)\leq\dim(A)+\dim(B\otimes_A k)~;
\end{equation}
en second lieu, d'après (16.2.6), on a
\begin{equation}
\label{0.17.3.3.4-fr}
\tag{17.3.3.4}
  \dim(A)\leq\operatorname{rg}_{k'}(V\otimes_k k').
\end{equation}
Enfin, dans l'anneau local $B\otimes_A k=B/\mathfrak{m}B$,
$\mathfrak{n}'=\mathfrak{n}/\mathfrak{m}B$ est l'idéal maximal, $k'$ le corps
résiduel et $\mathfrak{n}'/\mathfrak{n}'^2$ est isomorphe à
$\mathfrak{n}/(\mathfrak{n}^2+\mathfrak{m}B)
=W/\psi(V\otimes_k k')$~; on a donc, d'après (16.2.6)
\begin{equation}
\label{0.17.3.3.5-fr}
\tag{17.3.3.5}
  \dim(B\otimes_A k)
  \leq\operatorname{rg}_{k'}W
  -\operatorname{rg}_{k'}\psi(V\otimes_k k').
\end{equation}
Enfin, puisque $B$ est supposé régulier, on a
$\dim(B)=\operatorname{rg}_{k'}W$ (17.1.1)~; on conclut donc de (17.3.3.3),
(17.3.3.4) et (17.3.3.5) que l'on a
\[
  \operatorname{rg}_{k'}W
  \leq\dim(A)+\dim(B\otimes_A k)
  \leq\operatorname{rg}_{k'}W
  +\operatorname{rg}_{k'}(V\otimes_k k')
  -\operatorname{rg}_{k'}\psi(V\otimes_k k')
\]
et dire que les deux termes extrêmes de cette inégalité sont égaux signifie
que $\psi$ est \emph{injectif}. La condition a) entraîne donc nécessairement
que dans chacune des relations (17.3.3.3), (17.3.3.4) et (17.3.3.5), les deux
membres sont \emph{égaux}~; or, l'égalité dans (17.3.3.4) (resp. (17.3.3.5))
signifie que $A$ (resp. $B\otimes_A k$) est \emph{régulier} (17.1.1)~; on a
donc bien prouvé que a) entraîne e).
\end{proof}

\begin{proposition}[17.3.4]
\label{0.17.3.4-fr}
Soient $A$ un anneau local régulier de dimension $n$, $\mathfrak{m}$ son idéal
maximal. Pour tout $A$-module non nul de type fini $M$, on a
\begin{equation}
\label{0.17.3.4.1-fr}
\tag{17.3.4.1}
  \operatorname{prof}(M)+\operatorname{dim.proj}(M)=n.
\end{equation}
\end{proposition}

\begin{proof}
Raisonnons par récurrence sur $m=\operatorname{prof}(M)$. Si $m=0$, on sait
(16.4.6, (i)) qu'il existe un sous-$A$-module $N$ de $M$ isomorphe à
$k=A/\mathfrak{m}$~; appliquant la suite exacte des Tor à la suite exacte
$0\longrightarrow N\longrightarrow M\longrightarrow M/N\longrightarrow0$,
on obtient une suite exacte
\[
  \operatorname{Tor}^A_{n+1}(M/N,k)\longrightarrow
  \operatorname{Tor}^A_n(N,k)\longrightarrow
  \operatorname{Tor}^A_n(M,k),
\]
et l'hypothèse que $A$ est régulier entraîne
$\operatorname{Tor}^A_{n+1}(M/N,k)=0$ ((17.3.1.1) et (17.2.6)), donc
$\operatorname{Tor}^A_n(k,k)$ est isomorphe à un sous-$A$-module de
$\operatorname{Tor}^A_n(M,k)$~; appliquant de nouveau (17.3.1.1) et (17.2.6),
on voit que $\operatorname{Tor}^A_n(M,k)\neq0$, donc
$\operatorname{dim.proj}(M)\geq n$ (17.2.6)~; mais comme
$\operatorname{dim.proj}(M)\leq n$ par (17.3.1.1), on a bien
$\operatorname{dim.proj}(M)=n$.

Supposons maintenant $m>0$, et soit $x$ un élément $M$-régulier appartenant à
$\mathfrak{m}$~; on sait alors (16.4.6, (i)) que l'on a
$\operatorname{prof}(M/xM)=m-1$, et d'autre part (17.2.7)
$\operatorname{dim.proj}(M/xM)=\operatorname{dim.proj}(M)+1$~; l'hypothèse de
récurrence prouve aussitôt la relation (17.3.4.1).
\end{proof}

\oldpage[0\textsubscript{IV-1}]{50}

\begin{corollary}[17.3.5]
\label{0.17.3.5-fr}
\begin{enumerate}
  \item[\rm(i)] Soit $A$ un anneau local régulier de dimension $n$~; pour
  qu'un $A$-module $M\neq0$ de type fini soit libre, il faut et il suffit que
  $M$ soit un $A$-module de Cohen-Macaulay de dimension $n$.
  \item[\rm(ii)] Soient $A$ un anneau local régulier, $B$ un anneau local,
  $\rho:A\to B$ un homomorphisme local faisant de $B$ un $A$-module de type
  fini. Pour que $B$ soit un $A$-module libre, il faut et il suffit que $B$
  soit un anneau de Cohen-Macaulay et que $\rho$ soit injectif (ou, ce qui
  revient au même (16.1.5), que $\dim(B)=\dim(A)$).
\end{enumerate}
\end{corollary}

\begin{proof}
(i) Cela résulte de (17.3.4.1) et du fait que pour un $A$-module de type fini,
il revient au même de dire que ce module est projectif ou libre (Bourbaki,
\emph{Alg. comm.}, chap.~II, \S~3, n\textsuperscript{o}~2, cor.~2 de la
prop.~5)~; les $A$-modules libres de type fini $M$ sont donc caractérisés par
la relation $\operatorname{dim.proj}(M)=0$ (17.2.2, (i)).

(ii) Dire que $B$ est un anneau de Cohen-Macaulay équivaut à dire que $B$ est
un $A$-module de Cohen-Macaulay (16.5.3), donc il suffit d'appliquer (i),
puisque $\dim(B)=\dim(A)$ (16.1.5).
\end{proof}

\begin{env}[17.3.6]
\label{0.17.3.6-fr}
On dit qu'un anneau noethérien $A$ est \emph{régulier} si pour tout idéal
premier $\mathfrak{p}$ de $A$, l'anneau local $A_{\mathfrak{p}}$ est
régulier~; lorsque $A$ lui-même est un anneau local, il résulte de (17.3.2)
que cette définition est équivalente à celle de (17.1.2). Pour que $A$ soit
régulier, il suffit, en vertu de (17.3.2), que $A_{\mathfrak{m}}$ soit
régulier pour tout idéal maximal $\mathfrak{m}$ de $A$. En outre, il résulte
aussitôt de cette définition que pour toute partie multiplicative $S$ de
$A$, $S^{-1}A$ est régulier.
\end{env}

\begin{proposition}[17.3.7]
\label{0.17.3.7-fr}
Si $A$ est un anneau noethérien régulier, tout anneau de polynômes
$A[T_1,\ldots,T_n]$ est régulier.
\end{proposition}

\begin{proof}
Il suffit évidemment de prouver que l'anneau de polynômes $B=A[T]$ est
régulier~; comme $B$ est un $A$-module libre, pour tout idéal premier
$\mathfrak{q}$ de $B$, $B_{\mathfrak{q}}$ est un
$A_{\mathfrak{p}}$-module plat, où $\mathfrak{p}=\mathfrak{q}\cap A$
($0_{\mathrm I}$, 6.3.1)~; il suffit donc, en vertu de (17.1.10), de prouver
que $B_{\mathfrak{q}}/\mathfrak{p}B_{\mathfrak{q}}$ est régulier, et comme
cet anneau est un anneau local en un idéal premier de
$A_{\mathfrak{p}}[T]/\mathfrak{p}A_{\mathfrak{p}}[T]=k[T]$ (où $k$ est le
corps résiduel $A_{\mathfrak{p}}/\mathfrak{p}A_{\mathfrak{p}}$), il suffit
de prouver que $C=k[T]$ est régulier~; or, $C$ étant principal, les anneaux
locaux en les idéaux premiers de $C$ sont des anneaux de valuation discrète
ou un corps (pour l'idéal $(0)$), donc réguliers (17.1.4), ce qui achève la
démonstration.
\end{proof}

\begin{corollary}[17.3.8]
\label{0.17.3.8-fr}
Si $A$ est un anneau régulier, tout anneau de séries formelles
$A[[T_1,\ldots,T_n]]$ est régulier.
\end{corollary}

\begin{proof}
Soit $\mathfrak{J}$ l'idéal engendré par les $T_i$ dans l'anneau de polynômes
$A[T_1,\ldots,T_n]$~; comme ce dernier est régulier par (17.3.7), et que
$A[[T_1,\ldots,T_n]]$ est le complété de $A[T_1,\ldots,T_n]$ pour la
topologie $\mathfrak{J}$-préadique (Bourbaki, \emph{Alg. comm.}, chap.~III,
\S~2, n\textsuperscript{o}~12, \emph{Exemple 1}), la conclusion résulte du
\end{proof}

\begin{lemma}[17.3.8.1]
\label{0.17.3.8.1-fr}
Soient $A$ un anneau régulier, $\mathfrak{J}$ un idéal de $A$,
$\widehat{A}$ le séparé complété de $A$ pour la topologie
$\mathfrak{J}$-préadique. Alors $\widehat{A}$ est régulier.
\end{lemma}

\begin{proof}
Il suffit en effet (17.3.6) de voir que pour tout idéal maximal
$\mathfrak{n}$ de $\widehat{A}$, l'anneau local
$\widehat{A}_{\mathfrak{n}}$ est régulier~; on sait (Bourbaki,
\emph{Alg. comm.}, chap.~III, \S~3, n\textsuperscript{o}~4, prop.~8) que
$\mathfrak{n}$
\oldpage[0\textsubscript{IV-1}]{51}
est de la forme $\mathfrak{m}\widehat{A}$, où $\mathfrak{m}$ est un idéal
maximal de $A$ contenant $\mathfrak{J}$, et que l'homomorphisme canonique
$A\to\widehat{A}$ donne un homomorphisme \emph{injectif}
$A_{\mathfrak{m}}\to\widehat{A}_{\mathfrak{n}}$, tel que la topologie
$\mathfrak{n}\widehat{A}_{\mathfrak{n}}$-préadique de
$\widehat{A}_{\mathfrak{n}}$ induise sur $A_{\mathfrak{m}}$ la topologie
$\mathfrak{m}A_{\mathfrak{m}}$-préadique, et que $A_{\mathfrak{m}}$ soit
dense dans $\widehat{A}_{\mathfrak{n}}$ pour la topologie
$\mathfrak{n}\widehat{A}_{\mathfrak{n}}$-préadique~; par suite les
\emph{complétés} des anneaux locaux noethériens $A_{\mathfrak{m}}$ et
$\widehat{A}_{\mathfrak{n}}$ sont canoniquement \emph{isomorphes}. Or, par
hypothèse $A_{\mathfrak{m}}$ est régulier, donc il en est de même de son
complété (17.1.5) et comme le complété de $\widehat{A}_{\mathfrak{n}}$ est
régulier, il en est de même de $\widehat{A}_{\mathfrak{n}}$ par (17.1.5).
\end{proof}

\begin{corollary}[17.3.9]
\label{0.17.3.9-fr}
Soit $A$ un anneau noethérien, quotient d'un anneau noethérien régulier $B$.
Si $C$ est une $A$-algèbre de type fini, tout anneau de fractions $S^{-1}C$
de $C$ est quotient d'un anneau régulier.
\end{corollary}

\begin{proof}
En effet, $C$ est quotient d'un anneau de polynômes $A[T_1,\ldots,T_n]$~;
comme $A=B/\mathfrak{J}$, où $\mathfrak{J}$ est un idéal de $B$, on a
$A[T_1,\ldots,T_n]=B[T_1,\ldots,T_n]/\mathfrak{J}B[T_1,\ldots,T_n]$~; en
vertu de (17.3.7), on peut donc se borner au cas où $C$ est quotient d'un
anneau régulier $B'$~; mais si $S'$ est l'image réciproque de $S$ dans $B'$,
$S'$ est une partie multiplicative de $B'$ et $C$ est un anneau quotient de
$S'^{-1}B'$, si bien que finalement tout revient à voir que lorsque $C$ est
\emph{régulier}, il en est de même de $S^{-1}C$, ce qu'on a vu en (17.3.6).
\end{proof}

L'importance des quotients d'anneaux réguliers tient entre autres à la
propriété précédente et au fait que ce sont des anneaux \emph{caténaires}
(16.5.12).

Tous les anneaux importants dans les applications à la Géométrie algébrique
sont des quotients d'anneaux réguliers.
